The sound came across the water at dawn and Arya felt it in her teeth.

She had been at the shoreline since before first light, where she went most mornings now — down past the last of the camp’s picket lines to the shingle beach south of the harbor, because it was the one place within reach of the trenches where a person could stand and look at something other than a wall that would not open. She had watched Euron Greyjoy’s fleet ride at anchor there for four days, counting hulls and marking watches the way she counted and marked everything, and had begun, somewhere in the second day of it, to work out how a person might get aboard one of those ships if a person ever needed to.

Then the sound.

It was not a horn, whatever it had come out of. Arya had heard horns her whole life — hunting horns at Winterfell, watch-horns on the Wall’s memory, the flat brass of a Braavosi harbor calling ships in through fog — and none of them had ever gone \textit{through} her the way this went through her, a note that arrived in the bones of her jaw and the backs of her eyes rather than in her ears, and stayed there a long moment after it had stopped, ringing in a place that had no business ringing. Down the beach, a gull dropped out of the air mid-flight and did not get up again. Behind her, in the camp, she heard horses begin to scream.

She was already running before she had decided to.

\sceneBreak

Bran had gone rigid in his chair by the time she reached him, Sam holding his shoulders with both hands and the panic on his broad face plain enough to read from thirty feet, and Sansa arriving from the other direction with Tyrion hurrying behind her, and the whole camp coming awake into the particular disorder of people who had all just heard something they had no name for.

“It’s the horn,” Bran said, when he came back — and it took him longer than usual, long enough that Sam had begun saying his name over and over in the flat, frightened way people said names when they were afraid of being unheard. His eyes found Arya first, which she understood without needing it explained. “The Greyjoy. He’s sounded it. He’s reaching for Drogon.”

“Reaching,” Tyrion said. “Meaning what, exactly. Speak plainly, for once in your life—“

“Meaning Drogon isn’t answering Daenerys the way he should. Not entirely. Not anymore.” Bran’s voice was thin, and Arya thought — filing it the way she filed everything — that he sounded genuinely shaken for the first time since Winterfell. “They’re flying back from Dragonstone. He’s fighting it. She’s fighting it. I don’t know which of them is going to win and I don’t think either of them does either.”

Sansa’s face had gone to the flat stillness Arya had learned meant her sister was doing arithmetic rather than feeling anything about it yet. “Can it be broken?”

“The tales all say the same thing.” Sam, unexpectedly — and Arya watched him swallow, and understood he was reaching for the one thing he had that was useful. “Every version I ever read. The horn binds a dragon to whoever sounds it, and it costs the sounder terribly, and the binding lasts exactly as long as the man who blew it does. That’s the whole shape of it. It isn’t a spell laid on the dragon. It’s a leash held in a living hand.”

Nobody said anything for a moment. Arya felt every face around that chair arrive at the same conclusion at very nearly the same time, and found she had arrived there several seconds before any of them.

“He’s on his own ship,” she said. “\textit{Silence}. The black one with the red hull, three from the harbor mouth. He’s been aboard her every dawn for four days; I’ve watched him come up on deck.” She was already checking the knife at her hip, the small one, the one that mattered. “And there’s twenty-odd longships coming up the Blackwater under Greyjoy sail that aren’t his, which means in about an hour that harbor is going to be the most confused stretch of water in Westeros.”

“Arya.” Sansa’s voice, sharp. “You are not going to swim into an enemy fleet alone—“

“I’m going to walk,” Arya said. “Swimming’s for people who haven’t thought it through.”

\sceneBreak

The Ironborn kept a shore camp on the beach below the harbor, the way sailors always kept a shore camp when their captain had them anchored too long in one place — cookfires, a water detail, men rotating down off the ships to stretch legs that had spent four days on a deck. Arya had counted the rotation. She knew when the water detail came down and how long it took and how many men came with it, and she knew which of them walked back last and alone because he stopped to piss in the same stand of rocks every morning, a heavyset man with a black beard and a Greyjoy kraken sewn crooked on his jerkin.

He did not hear her. They rarely did.

She did it quickly, and cleanly, and with the small sharp motion the kindly man had spent a year drilling into her before she had understood what she was being drilled for, and afterward she knelt over the body in the cold rocks and did the other thing, the thing she had not done in a long while and had never once done without her hands shaking a little.

\textit{Valar dohaeris,} she thought, which was what you thought, when you did this, because someone had to think something and the man himself no longer could.

His face came away the way they always came away — easier than a person would imagine, which was the worst part of it — and when she rose and stood a moment in the rocks with the salt wind on skin that was not hers, she found she had to spend a small deliberate effort remembering the shape of the girl underneath it, the way she always had to now. There had been a time when that frightened her. She had made a kind of peace with it, the way she’d made peace with most of what Braavos had put into her, which was not the same as forgiving it.

He had been tongueless. She had known that before she chose him, had known it about all of them, because Euron Greyjoy’s crew were famous for it up and down every coast in Westeros — a captain who cut out his men’s tongues so they could never carry his secrets ashore. It was the cleverest thing about him and it was going to kill him, and Arya found, walking up the shingle in a dead man’s face with a dead man’s crooked kraken on her chest, that she appreciated the arrangement more than she’d expected to. No one would ask her to speak. No one aboard that ship \textit{could} ask her to speak. Euron had built himself a crew of men who could not be questioned, and had never once considered that a crew who could not be questioned also could not be verified.

She went up the gangplank with the last of the water detail and nobody looked at her twice.

\sceneBreak

\textit{Silence} did not creak. That was the first thing she noticed, standing on his deck, and the wrongness of it prickled at her worse than the horn had — a ship this size should have talked, timber and rope and the constant small complaint of a hull working against water, and this one made almost no sound at all beneath a full crew’s feet. Arya had spent enough of her life in enough terrible places to recognize when a man had built himself something to be admired rather than used, and to understand what such a man was likely to be like in the last few feet before he died.

The deck was in disorder, which helped. Yara Greyjoy’s longships were up the channel now, close enough to count, and Euron’s captains were shouting orders at men who could not shout back, and the whole harbor had begun to churn with the particular chaos of a fleet trying to decide whether it was about to fight. Arya moved through it the way she moved through everything, unhurried, carrying a water cask on one shoulder because a man carrying something was invisible in a way a man carrying nothing never was, and worked her way aft, and up, toward the raised deck at the stern where she had watched a one-eyed man stand every dawn for four days.

He was still there. He was not standing well.

That was the thing she had not expected, coming up the last few steps with the cask on her shoulder — that Euron Greyjoy, who she had watched hold that rail for four mornings like a man who owned the sea he was looking at, was gripping it now with both hands to stay upright. There was blood dried at the corners of his mouth and more of it fresh on his chin. His skin had gone a grey that Arya, who had seen a great many people at a great many stages of dying, recognized immediately. And there was a horn on the deck beside him, red-gold and black iron, and even lying still and silent on the boards it made the air near it feel wrong, the way meat left too long makes the air near it feel wrong.

He was, she understood with something that was not quite disappointment, already dying. She had come to kill a man who had spent himself down to nearly nothing an hour before she arrived.

She set the cask down.

“You,” Euron said, without turning. His voice had gone to wreckage. “Water. Here.”

She came up beside him. He held out one hand for a cup she did not have, and she took the hand instead — took the wrist, the way she’d been taught, thumb into the soft place where the tendons ran — and put the knife in under his ribs on the upward angle with her whole weight behind it, and held him against the rail while it happened so he could not fall and could not shout and could not do anything at all except turn his head and look at her.

His eye found her face. The dead crewman’s face. She watched him try to make sense of it — watched a man who had spent his entire life being the cleverest thing in every room he entered spend the last of himself on a puzzle he did not have time left to solve.

So she gave it to him. She had not planned to. It was not the discipline; it was the opposite of the discipline, and she would think about that later, at length. But she reached up with her free hand and took the dead man’s face off in front of him, and let Euron Greyjoy see exactly who was killing him.

“Arya Stark,” she said. “Of Winterfell. My brother died holding a gate you didn’t bother sailing to.”

Something moved across the ruin of his face that she would not have called fear. It was closer to recognition — the specific, exhausted recognition of a gambler shown a card he’d known all along was somewhere in the deck. His lips moved. She was near enough to hear it.

“\textit{Ah,}” he said. Just that. And then, with what he had left, something that might have been a laugh, or the beginning of one, and Arya understood in that moment that the fire she’d heard about, the one the horn lit inside a man, was still burning in him right up to the end and would go on burning for the fraction of a second longer than he did.

Then he was gone, and she let the weight of him down onto the deck, and turned around to find the entire stern of \textit{Silence} looking at her.

\sceneBreak

Two dozen men. Tongueless, every one. No shout had gone up because no shout could, and so the whole thing had happened in the middle of a crowded warship in almost perfect quiet, and now two dozen men who had served Euron Greyjoy stood in a loose ring around a girl holding a dead man’s face in one hand and a wet knife in the other, and none of them made a sound, and none of them moved.

Arya had been in a great many rooms she was not certain she would leave. She had learned that the thing to do in such a room was to be the person who spoke first.

“He’s dead,” she said, to all of them, in her own voice, which she had not used aloud in an hour. “The horn’s finished. Whatever he had a hold of, he hasn’t got it anymore.” She let the face fall to the deck, because holding it was doing her no favors with this audience. “Your captain cut your tongues out so you couldn’t tell anyone what he was. I don’t expect any of you loved him for it.”

Nobody moved. Out across the water, close now, close enough that Arya could hear the oars, the first of the Greyjoy longships was coming in past the harbor mouth under a banner that was not Euron’s, and every man on that deck turned his head toward it, and Arya understood — watching two dozen silent faces do the same arithmetic at the same time — that she had not actually saved herself at all. She had only made the question interesting, and someone else was about to arrive to answer it.

She picked her spot at the rail where she could get into the water fast if the answer went badly, and waited, and found that her hands had begun to shake at last, now that there was nothing left to do with them.

Somewhere far above and behind her, faint, ragged, wrong — but there, and getting closer — she heard a dragon scream.